\chapter{Configuration and environments}
\label{ch:configuration}

\section{Environments}

Yap has one backend environment: a hosted Supabase project, which every build of the app uses
(\cref{fig:deployment}). The app reaches it through the URL and publishable key in
\code{src/config/supabase.ts}. There is no separate staging project; behaviour that must be proved
against a database is proved against a disposable local Supabase stack (\cref{ch:testing}).
\Cref{app:config} lists every configuration value.

\begin{figure}[ht]
  \centering
  % Environments and how code and schema reach them (Figure: deployment).
\begin{tikzpicture}[
    env/.style={box, font=\sffamily\scriptsize, align=flush left, inner sep=6pt, anchor=north},
    lab/.style={font=\sffamily\scriptsize, text=ink, fill=white, inner sep=1.2pt, align=center},
  ]
  \node[env, clientbox, text width=42mm] (dev) at (2.4,0)
    {\textbf{\footnotesize Developer machine}\\[3pt]
     repository: \texttt{src}, \texttt{supabase/migrations}, \texttt{scripts}\\
     Node 22, \texttt{npm install}\\
     \texttt{npx expo start}: development server\\
     \texttt{npm run icons}, \texttt{npm run mascot}};
  \node[env, clientbox, text width=46mm] (phones) at (12.3,0)
    {\textbf{\footnotesize Phones}\\[3pt]
     Expo Go on iPhone and Android\\
     app bundle from the development server\\
     session kept in AsyncStorage};
  \node[env, serverbox, text width=46mm] (hosted) at (12.3,-3.0)
    {\textbf{\footnotesize Hosted Supabase project}\\[3pt]
     Auth: email and password, 8-digit codes\\
     \quad custom SMTP: Resend, \texttt{mail.samarxkx.dev}\\
     Postgres with PostGIS and \texttt{pg\_cron}\\
     Realtime; Storage bucket \texttt{photos}\\
     dashboard: Auth settings, email templates};
  \node[env, text width=46mm, dashed] (local) at (12.3,-6.6)
    {\textbf{\footnotesize Local Supabase stack (Docker)}\\[3pt]
     created for one change, then torn down\\
     every migration, plus a stand-in for the\\
     \quad platform helper \texttt{rls\_auto\_enable}\\
     used for behaviour proofs only};
  \node[env, text width=42mm] (gh) at (2.4,-4.4)
    {\textbf{\footnotesize GitHub repository}\\[3pt]
     pull requests through the review pipeline\\
     no continuous-integration workflows};

  \draw[flow] (dev.east |- phones.west) -- node[lab, above=2pt] {app bundle over the local network} (phones.west);
  \draw[flow, <->] (phones.south) -- node[lab, right=2pt] {HTTPS and WebSocket} (hosted.north);
  \draw[flow] (dev.east |- 0,-1.75) -- (7.2,-1.75) |- (hosted.west |- 0,-3.9);
  \node[lab, anchor=south west, align=left, fill=none] at (7.25,-3.85) {migrations, not\\applied with\\\texttt{supabase db push}};
  \draw[softflow] (dev.east |- 0,-2.15) -- (6.6,-2.15) |- (local.west |- 0,-7.5);
  \node[lab, anchor=south west, align=left] at (6.7,-7.45) {migrations, for proofs};
  \draw[flow, <->] (dev.south) -- node[lab, left=2pt] {push; pull\\requests} (gh.north);
\end{tikzpicture}

  \caption{Environments. Phones run the app in Expo Go against the hosted project; migrations reach the
    hosted project from the repository; a throwaway local stack is used only to prove changes.}
  \label{fig:deployment}
\end{figure}

\section{Supabase dashboard settings}
\label{sec:dashboard}

Some configuration lives in the Supabase dashboard rather than in the repository. The README records the
settings that password sign-in needs, and the project's email delivery is configured alongside them
(\cref{tab:dashboard}).

\begin{table}[ht]
  \centering
  \caption{Settings made in the Supabase dashboard.}
  \label{tab:dashboard}
  \tablefont
  \begin{tabularx}{\textwidth}{@{}L{50mm}Y@{}}
    \toprule
    \textsf{\textbf{Setting}} & \textsf{\textbf{Value}} \\
    \midrule
    Auth, email provider & Enabled for password sign-in; \ui{Confirm email} on; minimum password length 8 \\
    Template: confirm sign-up & Shows the 8-digit code (\texttt{\{\{ .Token \}\}}) \\
    Template: reset password & Shows the 8-digit code instead of a link, with a subject such as \ui{Reset your Yap password} \\
    Template: magic link & No longer used by the app \\
    Email delivery & Custom SMTP through Resend, sending from \code{mail.samarxkx.dev} \\
    Apple and Google providers & Not configured; the app's buttons are behind a flag \\
    \bottomrule
  \end{tabularx}
\end{table}

The \code{pg_cron} extension and its schedule, the \code{photos} bucket, the Realtime publication and all
policies are created by the migrations, not by hand.

\section{Database migrations}
\label{sec:migrations}

The schema is defined by eight SQL migrations in \code{supabase/migrations}, listed in
\cref{app:migrations}. Each later migration replaces functions with \code{CREATE OR REPLACE}, which keeps
their grants, or drops and recreates a function whose arguments or return type change and then grants it
again in the same way as the original.

Two practices govern how migrations reach the hosted project. First, the hosted project's recorded
migration versions do not match the file names in the repository, so a new migration is applied in the
same way as the earlier ones rather than with \code{supabase db push}. Second, an app change and the
migration it depends on ship together: after the server-rules migration, for example, \code{join_event}
takes the phone's position, so an older app build can no longer join events.

\code{src/data/database.types.ts} holds TypeScript types generated from the schema, which give the
Supabase client typed tables and function arguments. It is regenerated with
\code{supabase gen types typescript} (or the equivalent Supabase tool) whenever a migration changes tables
or functions.

\section{Local development}

The README describes the development loop:

\begin{lstlisting}[language=bash, caption={Running the app for development.}, label={lst:run}]
nvm use            # Node 22, from .nvmrc
npm install
npx expo start     # scan the QR code with Expo Go on an iPhone or Android phone
\end{lstlisting}

The app runs in Expo Go against the hosted project. Exercising the core loop needs two phones signed in
with different email addresses, within about 100 metres of each other, both live. The native
\code{ios} and \code{android} folders are generated and not committed.

A local Supabase stack can run the migrations for testing, with one adjustment: migration
\code{20261004210004_lock_platform_helper.sql} revokes \code{public.rls_auto_enable()}, a function the
hosted platform installs, so a local stack needs a stand-in for it.

\section{App configuration}

\code{app.json} configures the build (\cref{tab:app-json}).

\begin{table}[ht]
  \centering
  \caption{App configuration in \code{app.json}.}
  \label{tab:app-json}
  \tablefont
  \begin{tabularx}{\textwidth}{@{}L{44mm}Y@{}}
    \toprule
    \textsf{\textbf{Key}} & \textsf{\textbf{Value}} \\
    \midrule
    Name, slug, version & \code{Yap}, \code{yap}, \code{1.0.0}; portrait only; URL scheme \code{yap} \\
    Appearance & \code{automatic}: light or dark from the phone's setting \\
    iOS icon & The \code{assets/expo.icon} bundle, which draws the brand mark over a green fill \\
    Android icon & Adaptive icon with foreground, background and monochrome layers; background \texttt{\#15803D} \\
    Splash screen & Brand mark on \texttt{\#FAF8F5} (light) or \texttt{\#181715} (dark) \\
    Location permission & When in use only, with the text ``Yap uses your location only while you are available, to show you people who are up for a chat in the same place. Nobody sees where you are.'' \\
    Photo and camera permissions & For choosing a profile photo, and for taking one and a verification selfie \\
    Plugins & Expo Router, splash screen, location, image picker, date and time picker \\
    Experiments & Typed routes; React Compiler \\
    \bottomrule
  \end{tabularx}
\end{table}

\section{Feature flags}

Two feature flags in \code{src/config/features.ts} are both off: \code{selfieVerification} and
\code{socialSignIn} (\cref{sec:feature-flags}). Turning a flag on shows the corresponding screens and
buttons; neither feature has a working backend yet.

\section{Asset pipelines and scripts}

Two scripts regenerate assets from source drawings, so the drawings are the single source of truth
(\cref{tab:scripts}).

\begin{description}
  \item[Icons.] \code{npm run icons} runs \code{scripts/make-icons.sh}, which uses ImageMagick 7 to draw
    the brand mark in \code{assets/brand/yap-mark.svg} into the app icon, the three Android adaptive-icon
    layers, the light and dark splash images and the favicon, and copies the SVG into the iOS icon bundle.
    Colours come from the theme tokens. It strips metadata from every image, so re-running it without
    changes leaves the working tree clean.
  \item[Mascot.] \code{npm run mascot} runs \code{scripts/make-mascot.mjs}, which reads each drawing in
    \code{assets/mascot}, removes titles and comments, and writes them as strings into
    \code{src/components/mascot/mascot-art.ts}. Each file name becomes a pose, drawn with
    \code{react-native-svg} without any extra bundler setup.
\end{description}

\begin{table}[ht]
  \centering
  \caption{Package scripts.}
  \label{tab:scripts}
  \tablefont
  \begin{tabularx}{\textwidth}{@{}L{34mm}L{50mm}Y@{}}
    \toprule
    \textsf{\textbf{Script}} & \textsf{\textbf{Command}} & \textsf{\textbf{Purpose}} \\
    \midrule
    \code{start} & \code{expo start} & Development server \\
    \code{ios}, \code{android} & \code{expo start --ios}, \code{--android} & Development server opening a simulator or device \\
    \code{web} & \code{expo start --web} & From the Expo template; the web target is not supported (\cref{ch:limitations}) \\
    \code{lint} & \code{expo lint} & ESLint with the Expo configuration \\
    \code{icons} & \code{sh scripts/make-icons.sh} & Regenerate icons, splash images and favicon \\
    \code{mascot} & \code{node scripts/make-mascot.mjs} & Regenerate the mascot drawings module \\
    \bottomrule
  \end{tabularx}
\end{table}
